\section{Mathematical methods}
\label{app:math_methods}

This appendix provides full mathematical specifications of the algorithms summarized in section 2 of the main text. Each subsection corresponds to one component of the periodic and rhythmic discharge characterization pipeline, including network architectures, evidence-combination formulas, dynamic programming, frequency estimation, spatial localization, and BIPD detection. Equation numbering is local to each subsection and follows the source method writeups.

\subsection{ChannelPD-Net: per-channel CNN with temporal attention}
\label{app:channelpdnet}

\subsubsection{Overview}

ChannelPD-Net is a lightweight 1D convolutional neural network with temporal attention pooling that operates independently on each bipolar EEG channel. Given a single 10-second z-scored channel, it jointly predicts (i) the probability that a periodic discharge pattern is present and (ii) the log-frequency of that pattern. The model has approximately 50,000 trainable parameters and is deployed as a 5-fold patient-stratified ensemble. Its per-channel outputs serve as the foundation for lateralization, spatial localization, and frequency estimation in the full characterization pipeline.

\subsubsection{Input representation}

Each input is a single bipolar EEG channel resampled to $F_s = 200$ Hz, yielding $N = 2000$ samples for a 10-second epoch. The channel is z-scored to zero mean and unit variance:
\[ x_n = \frac{x_n^{\mathrm{raw}} - \hat{\mu}}{\hat{\sigma}}, \quad n = 1, \ldots, N \tag{1} \]
where $\hat{\mu}$ and $\hat{\sigma}$ are the sample mean and standard deviation of the channel. The model receives $\mathbf{x} \in \mathbb{R}^{1 \times N}$ as a single-channel 1D signal.

\subsubsection{Architecture}

\paragraph{Convolutional backbone.}
The backbone consists of four 1D convolutional blocks, each applying convolution, batch normalization, GELU activation, and dropout. Stride-2 convolutions progressively downsample the temporal dimension by a factor of $2^4 = 16$.
\[ \mathbf{h}^{(1)} = \mathrm{Drop}_{0.1}\!\Big(\mathrm{GELU}\!\big(\mathrm{BN}\!\big(\mathrm{Conv1d}(\mathbf{x};\; C_{\mathrm{out}}{=}16,\; k{=}51,\; s{=}2,\; p{=}25)\big)\big)\Big) \tag{2} \]
\[ \mathbf{h}^{(2)} = \mathrm{Drop}_{0.1}\!\Big(\mathrm{GELU}\!\big(\mathrm{BN}\!\big(\mathrm{Conv1d}(\mathbf{h}^{(1)};\; 32,\; k{=}25,\; s{=}2,\; p{=}12)\big)\big)\Big) \tag{3} \]
\[ \mathbf{h}^{(3)} = \mathrm{Drop}_{0.1}\!\Big(\mathrm{GELU}\!\big(\mathrm{BN}\!\big(\mathrm{Conv1d}(\mathbf{h}^{(2)};\; 64,\; k{=}13,\; s{=}2,\; p{=}6)\big)\big)\Big) \tag{4} \]
\[ \mathbf{h}^{(4)} = \mathrm{Drop}_{0.2}\!\Big(\mathrm{GELU}\!\big(\mathrm{BN}\!\big(\mathrm{Conv1d}(\mathbf{h}^{(3)};\; 64,\; k{=}7,\; s{=}2,\; p{=}3)\big)\big)\Big) \tag{5} \]
After the backbone, $\mathbf{h}^{(4)} \in \mathbb{R}^{64 \times T'}$ where $T' = \lceil N / 16 \rceil = 125$.

\paragraph{Temporal attention pooling.}
Rather than global average pooling, we employ a learned temporal attention mechanism that allows the model to focus on the most informative time steps within the 10-second window. A $1 \times 1$ convolution projects each 64-dimensional feature vector to a scalar logit, followed by softmax normalization across time:
\[ \alpha_t = \frac{\exp\!\big(w^{\top} \mathbf{h}^{(4)}_t + b\big)}{\sum_{t'=1}^{T'} \exp\!\big(w^{\top} \mathbf{h}^{(4)}_{t'} + b\big)}, \quad t = 1, \ldots, T' \tag{6} \]
where $w \in \mathbb{R}^{64}$ and $b \in \mathbb{R}$ are the parameters of $\mathrm{Conv1d}(64 \to 1, k{=}1)$. The context vector is the attention-weighted sum of feature vectors:
\[ \mathbf{z} = \sum_{t=1}^{T'} \alpha_t \, \mathbf{h}^{(4)}_t \in \mathbb{R}^{64} \tag{7} \]

\paragraph{Dual-head prediction.}
Two independent linear heads operate on the context vector $\mathbf{z}$.

\textbf{PD detection head.} Produces a probability that the channel contains a periodic discharge pattern:
\[ p_{\mathrm{pd}} = \sigma\!\big(\mathbf{w}_{\mathrm{pd}}^{\top} \mathbf{z} + b_{\mathrm{pd}}\big) \in [0, 1] \tag{8} \]

\textbf{Frequency estimation head.} Predicts the log-frequency of the discharge pattern (trained only on PD-positive channels):
\[ \hat{f}_{\log} = \mathbf{w}_f^{\top} \mathbf{z} + b_f \in \mathbb{R} \tag{9} \]
where $\hat{f}_{\log} \approx \log(f_{\mathrm{gt}})$ in Hz.

\subsubsection{Training}

\paragraph{Loss function.}
The multi-task loss jointly optimizes detection and frequency estimation:
\[ \mathcal{L} = \mathrm{BCE}(p_{\mathrm{pd}},\; y_{\mathrm{pd}}) + \lambda \cdot m \cdot \mathrm{MSE}(\hat{f}_{\log},\; \log f_{\mathrm{gt}}) \tag{10} \]
where $y_{\mathrm{pd}} \in \{0, 1\}$ is the binary PD label, $f_{\mathrm{gt}}$ is the expert-annotated frequency, $m \in \{0, 1\}$ is a mask that is 1 only when $y_{\mathrm{pd}} = 1$ and a valid frequency label exists, and $\lambda = 0.5$ balances the two tasks.

\paragraph{Cross-validation and ensembling.}
Training uses 5-fold patient-stratified cross-validation, ensuring no patient appears in both training and validation folds. At inference, the ensemble prediction is the arithmetic mean across all five fold models:
\[ \bar{p}_{\mathrm{pd}} = \frac{1}{5} \sum_{k=1}^{5} p_{\mathrm{pd}}^{(k)}, \qquad \bar{f}_{\log} = \frac{1}{5} \sum_{k=1}^{5} \hat{f}_{\log}^{(k)} \tag{11} \]

\paragraph{Data augmentation.}
During training, three augmentations are applied stochastically:
\begin{itemize}
  \item \textbf{Amplitude scaling:} $x \leftarrow s \cdot x$ with $s \sim \mathcal{U}(0.8, 1.2)$
  \item \textbf{Additive Gaussian noise:} SNR drawn from $\mathcal{U}(20, 40)$ dB
  \item \textbf{Circular time shift:} $x \leftarrow \mathrm{roll}(x, \delta)$ with $\delta \sim \mathcal{U}(-50, 50)$ samples
\end{itemize}

\subsubsection{Downstream usage}

\paragraph{Laterality detection.}
For LPD segments, laterality is determined by comparing the mean PD probability of left-hemisphere channels versus right-hemisphere channels:
\[ \mathrm{laterality} = \begin{cases} \texttt{left} & \text{if } \bar{p}_{\mathrm{left}} > \bar{p}_{\mathrm{right}} \\ \texttt{right} & \text{otherwise} \end{cases} \tag{12} \]
where $\bar{p}_{\mathrm{left}} = \frac{1}{|\mathcal{L}|}\sum_{c \in \mathcal{L}} p_{\mathrm{pd}}^{(c)}$ and $\mathcal{L}, \mathcal{R}$ are the index sets of left and right hemisphere channels in the 18-channel bipolar montage.

\paragraph{CNN frequency estimate.}
The patient-level CNN frequency estimate is a PD-weighted log-space average across all channels:
\[ f_{\mathrm{CNN}} = \exp\!\left(\frac{\sum_{c=1}^{C} p_{\mathrm{pd}}^{(c)} \cdot \hat{f}_{\log}^{(c)}}{\sum_{c=1}^{C} p_{\mathrm{pd}}^{(c)}}\right) \tag{13} \]
This weighting ensures that channels with high PD probability contribute more to the frequency estimate, effectively ignoring channels dominated by artifact or background activity. The final estimate is clipped to the physiological range $[0.3, 3.5]$ Hz.

\subsubsection{Performance}

ChannelPD-Net achieves a channel-level PD detection AUC of approximately 0.82 and a patient-level AUC of approximately 0.93. The CNN frequency estimate alone achieves a Spearman correlation of approximately 0.55 with expert frequency ratings; this improves substantially when combined with the ACF frequency estimate (see CNN+ACF Frequency Ensemble method).

\subsection{HemiCET-UNet: hemisphere-restricted convolutional evidence-trace U-Net}
\label{app:hemicet}

\subsubsection{Overview}

HemiCET-UNet is a 1D U-Net encoder-decoder architecture that produces a dense, frame-level discharge evidence trace from a single z-scored EEG channel. Unlike ChannelPD-Net, which provides a single summary probability per channel, HemiCET-UNet outputs a temporal signal $E(t) \in [0,1]^N$ at the full input resolution ($N = 2000$ samples, 200 Hz), indicating the instantaneous likelihood of a periodic discharge at each time point. The model has approximately 525,000 trainable parameters and is deployed as a 5-fold patient-stratified ensemble.

\subsubsection{Input representation}

The input is identical to ChannelPD-Net: a single bipolar EEG channel $\mathbf{x} \in \mathbb{R}^{1 \times N}$ with $N = 2000$ samples at $F_s = 200$ Hz, z-scored per-channel:
\[ x_n = \frac{x_n^{\mathrm{raw}} - \hat{\mu}}{\hat{\sigma}}, \quad n = 1, \ldots, N \tag{1} \]

\subsubsection{Architecture}

\paragraph{Encoder.}
The encoder mirrors the ChannelPD-Net backbone: four convolutional stages with stride-2 downsampling, batch normalization, and GELU activation. Each stage halves the temporal dimension; see table~\ref{tab:hemicet_encoder}.

\begin{table}[h]
\centering
\caption{HemiCET-UNet encoder stages.}
\label{tab:hemicet_encoder}
\begin{tabular}{lll}
\hline
Stage & Operation & Output shape \\
\hline
$\mathbf{e}_1$ & $\mathrm{Conv1d}(1 \to 16,\; k{=}51,\; s{=}2,\; p{=}25) \to \mathrm{BN} \to \mathrm{GELU}$ & $16 \times 1000$ \\
$\mathbf{e}_2$ & $\mathrm{Conv1d}(16 \to 32,\; k{=}25,\; s{=}2,\; p{=}12) \to \mathrm{BN} \to \mathrm{GELU}$ & $32 \times 500$ \\
$\mathbf{e}_3$ & $\mathrm{Conv1d}(32 \to 64,\; k{=}13,\; s{=}2,\; p{=}6) \to \mathrm{BN} \to \mathrm{GELU}$ & $64 \times 250$ \\
$\mathbf{e}_4$ & $\mathrm{Conv1d}(64 \to 64,\; k{=}7,\; s{=}2,\; p{=}3) \to \mathrm{BN} \to \mathrm{GELU}$ & $64 \times 125$ \\
\hline
\end{tabular}
\end{table}

\paragraph{Decoder with skip connections.}
The decoder progressively upsamples back to the input resolution using transposed convolutions. At each stage, the upsampled feature map is concatenated with the corresponding encoder feature map (skip connection) and fused through a $1 \times 3$ convolution; see table~\ref{tab:hemicet_decoder}.

\begin{table}[h]
\centering
\caption{HemiCET-UNet decoder stages with skip connections. ConvT $=$ ConvTranspose1d. All ConvT layers use kernel $k{=}4$, stride $s{=}2$, padding $p{=}1$. All fusion convolutions are followed by batch normalization and GELU activation; the final $1 \times 1$ convolution is followed by a sigmoid.}
\label{tab:hemicet_decoder}
\begin{tabular}{lllll}
\hline
Stage & Upsample & Skip concat & Fusion conv & Output shape \\
\hline
$\mathbf{d}_4$ & ConvT$(64 \to 64)$ & $[\mathbf{d}_4 \,\|\, \mathbf{e}_3]$ (128 ch) & Conv1d$(128 \to 64,\, k{=}3)$  & $64 \times 250$ \\
$\mathbf{d}_3$ & ConvT$(64 \to 32)$ & $[\mathbf{d}_3 \,\|\, \mathbf{e}_2]$ (64 ch)  & Conv1d$(64 \to 32,\, k{=}3)$   & $32 \times 500$ \\
$\mathbf{d}_2$ & ConvT$(32 \to 16)$ & $[\mathbf{d}_2 \,\|\, \mathbf{e}_1]$ (32 ch)  & Conv1d$(32 \to 16,\, k{=}3)$   & $16 \times 1000$ \\
$\mathbf{d}_1$ & ConvT$(16 \to 8)$  & ---                                            & Conv1d$(8 \to 1,\, k{=}1)$     & $1 \times 2000$ \\
\hline
\end{tabular}
\end{table}

When encoder and decoder feature maps differ in length due to rounding, the longer map is truncated to match.

\paragraph{Output.}
The final layer applies a sigmoid activation to produce the evidence trace:
\[ E(t) = \sigma\!\big(\mathrm{Conv1d}(8 \to 1,\; k{=}1)\big) \in [0, 1]^N \tag{2} \]

\subsubsection{Training}

\paragraph{Target construction.}
Training targets are constructed from expert-annotated discharge times $\{t_1^*, t_2^*, \ldots, t_K^*\}$ (in samples). Each discharge generates a Gaussian peak in the target trace:
\[ y(t) = \max_{k=1}^{K} \exp\!\left(-\frac{(t - t_k^*)^2}{2\sigma_{\mathrm{target}}^2}\right) \tag{3} \]
where $\sigma_{\mathrm{target}} = 2$ samples (10 ms at 200 Hz), producing narrow, well-localized targets.

\paragraph{Loss function.}
The composite loss has three terms:
\[ \mathcal{L} = \mathcal{L}_{\mathrm{BCE}} + \gamma_1 \cdot \mathcal{L}_{\mathrm{sharp}} + \gamma_2 \cdot \mathcal{L}_{\mathrm{floor}} \tag{4} \]

\textbf{Term 1: Weighted binary cross-entropy.} Because discharge events are temporally sparse (a few dozen samples out of 2000), we apply asymmetric weighting with positive weight $w_+ = 20$:
\[ \mathcal{L}_{\mathrm{BCE}} = -\frac{1}{N} \sum_{t=1}^{N} \Big[ w_+ \cdot y(t) \log E(t) + (1 - y(t)) \log(1 - E(t)) \Big] \tag{5} \]

\textbf{Term 2: Sharpness penalty.} Encourages sparse evidence by penalizing the mean activation:
\[ \mathcal{L}_{\mathrm{sharp}} = \frac{1}{N} \sum_{t=1}^{N} E(t) \tag{6} \]
with $\gamma_1 = 0.1$.

\textbf{Term 3: Floor loss.} At expert-labeled discharge locations, the CET evidence should be at least as strong as the handcrafted HPP evidence, ensuring the learned model does not miss obvious discharges:
\[ \mathcal{L}_{\mathrm{floor}} = \frac{1}{K} \sum_{k=1}^{K} \max\!\big(0,\; E_{\mathrm{hpp}}(t_k^*) - E(t_k^*)\big)^2 \tag{7} \]
with $\gamma_2 = 0.05$.

\paragraph{Data augmentation.}
Four augmentation strategies are applied stochastically during training:
\begin{enumerate}
  \item \textbf{Amplitude scaling:} $x \leftarrow s \cdot x$ with $s \sim \mathcal{U}(0.7, 1.3)$
  \item \textbf{Additive Gaussian noise:} $x \leftarrow x + \epsilon$ with $\epsilon \sim \mathcal{N}(0, 0.05^2)$
  \item \textbf{Channel dropout:} each channel is zeroed with probability $p = 0.15$
  \item \textbf{Discharge jitter:} expert discharge times are perturbed by $\delta \sim \mathcal{N}(0, \sigma_j^2)$ with $\sigma_j = 1$ sample (5 ms), modeling annotation uncertainty
\end{enumerate}

\paragraph{Cross-validation and ensembling.}
The model is trained with 5-fold patient-stratified cross-validation. At inference, the ensemble evidence trace is the arithmetic mean:
\[ E_{\mathrm{ensemble}}(t) = \frac{1}{5} \sum_{k=1}^{5} E^{(k)}(t) \tag{8} \]

\subsubsection{Inference pipeline}

\begin{verbatim}
Input: single channel x in R^{2000}
  1. z-score normalize x
  2. For each fold model k = 1, ..., 5:
       E^(k) = CETUNet_k(x)           // forward pass
  3. E_ensemble = mean(E^(1), ..., E^(5))
Output: E_ensemble in [0,1]^{2000}    // frame-level evidence
\end{verbatim}

The per-channel evidence traces are subsequently aggregated across channels and combined with handcrafted evidence (see Dynamic Programming method) for discharge detection.

\subsubsection{Performance}

HemiCET-UNet evidence, when combined with handcrafted HPP evidence via the product-boost formula and processed through dynamic programming, achieves a discharge detection F1 of approximately 0.74 at a tolerance of $\pm 100$ ms. Using CET evidence alone (without HPP) yields F1 approximately 0.65, demonstrating that the learned and handcrafted evidence sources provide complementary information.

\subsection{Handcrafted Peak Prior (HPP) and product-boosted evidence combination}
\label{app:evidence_combination}

\subsubsection{Overview}

We combine two complementary evidence signals for discharge detection: a handcrafted heuristic peak-pointiness (HPP) trace and a learned convolutional evidence trace (CET) from the HemiCET-UNet ensemble. HPP excels at detecting sharp electrographic transients via morphological features, while CET captures complex learned discharge patterns. The two are fused using a product-boosted maximum operator, which selects the stronger signal at each timepoint while amplifying regions where both methods agree via a quadratic interaction term. This combination yields a single evidence trace suitable for downstream dynamic programming (DP) path optimization.

\subsubsection{Input}

\begin{itemize}
  \item Single-hemisphere EEG (8 bipolar channels), 10 s at 200 Hz
\end{itemize}

\subsubsection{Mathematical formulation}

\paragraph{HPP evidence: pointiness trace.}
For each local maximum at time $t_p$ in the 20 Hz lowpass-filtered signal, compute the pointiness score from its prominence $\rho(t_p)$ and half-prominence width $w(t_p)$:
\[ p(t_p) = \frac{\rho(t_p)^2}{w(t_p)} \tag{1} \]
The pointiness trace $P(t)$ is constructed by placing $p(t_p)$ at each peak location and zero elsewhere, then taking the channel-wise maximum across all 8 channels.

\paragraph{Teager-Kaiser energy operator (TKEO).}
The TKEO captures instantaneous energy and is sensitive to rapid amplitude and frequency changes. Applied to the 20 Hz lowpass signal $x(t)$:
\[ \mathrm{tkeo}[n] = \left|x[n]^2 - x[n-1] \cdot x[n+1]\right| \tag{2} \]
The TKEO trace is the channel-wise maximum across all 8 channels.

\paragraph{HPP combination.}
The two components are z-scored and combined as a weighted sum:
\[ E_{\mathrm{hpp}}(t) = 0.6 \cdot \tilde{P}(t) + 0.4 \cdot \widetilde{\mathrm{tkeo}}(t) \tag{3} \]
where $\tilde{\cdot}$ denotes z-scoring (zero mean, unit variance). The combined trace is smoothed with a Gaussian kernel ($\sigma = 3$ samples, i.e., 15 ms at 200 Hz):
\[ E_{\mathrm{hpp}}(t) \leftarrow (E_{\mathrm{hpp}} * g_\sigma)(t), \quad g_\sigma(t) = \frac{1}{\sigma\sqrt{2\pi}}\exp\!\left(-\frac{t^2}{2\sigma^2}\right) \tag{4} \]

\paragraph{CET evidence.}
The CET evidence $E_{\mathrm{cet}}(t)$ is the output of the HemiCET-UNet ensemble, a U-Net trained to predict discharge probability at each timepoint. The ensemble averages predictions from multiple folds for robustness.

\paragraph{Preprocessing.}
Both evidence traces are normalized to $[0,1]$:
\[ E_{\mathrm{hpp}}(t) \leftarrow \frac{E_{\mathrm{hpp}}(t) - \min(E_{\mathrm{hpp}})}{\max(E_{\mathrm{hpp}}) - \min(E_{\mathrm{hpp}})} \tag{5} \]
\[ E_{\mathrm{cet}}(t) \leftarrow \frac{E_{\mathrm{cet}}(t) - \min(E_{\mathrm{cet}})}{\max(E_{\mathrm{cet}}) - \min(E_{\mathrm{cet}})} \tag{6} \]
The CET trace is then thresholded at the 80th percentile of its non-zero values to suppress low-confidence predictions, and a hard floor is applied:
\[ \tau_{80} = \mathrm{percentile}_{80}\!\left(\{E_{\mathrm{cet}}(t) : E_{\mathrm{cet}}(t) > 0\}\right) \tag{7} \]
\[ E_{\mathrm{cet}}(t) \leftarrow \begin{cases} E_{\mathrm{cet}}(t) & \text{if } E_{\mathrm{cet}}(t) \geq \max(\tau_{80},\; 0.3) \\ 0 & \text{otherwise} \end{cases} \tag{8} \]

\paragraph{Product-boosted max combination.}
The final evidence trace is:
\[ E(t) = \max\!\left(E_{\mathrm{hpp}}(t),\; E_{\mathrm{cet}}(t)\right) + \lambda \cdot E_{\mathrm{hpp}}(t) \cdot E_{\mathrm{cet}}(t) \tag{9} \]
where $\lambda = 3.0$ is an empirically optimized boost coefficient.

The two terms serve distinct roles:
\begin{enumerate}
  \item \textbf{Max term}: At each timepoint, the stronger evidence source is selected. This ensures that a discharge detected by either method alone is preserved in the combined trace, preventing loss of sensitivity.
  \item \textbf{Product term}: The multiplicative interaction $E_{\mathrm{hpp}}(t) \cdot E_{\mathrm{cet}}(t)$ is nonzero only when both methods provide positive evidence at the same timepoint. Since both inputs are in $[0,1]$, the product is also in $[0,1]$ and is large only when both are simultaneously large. The coefficient $\lambda = 3.0$ amplifies this agreement signal, creating a strong synergistic boost.
\end{enumerate}

The product term effectively implements a soft AND gate: timepoints where both HPP and CET agree receive substantially higher evidence than those supported by only one method. This is motivated by the observation that true discharges activate both the morphological (sharp transient) and learned (contextual pattern) detectors, while artifacts tend to trigger only one.

\subsubsection{Algorithm}

\begin{verbatim}
EVIDENCE_COMBINATION(eeg[c,t]):
    # --- HPP Evidence ---
    1. x_lp = lowpass(eeg, 20 Hz)
    2. For each channel c:
         Find local maxima {t_p}
         Compute pointiness: p(t_p) = prominence(t_p)^2 / width(t_p)     [Eq. (1)]
         Compute TKEO: tkeo[n] = |x[n]^2 - x[n-1]*x[n+1]|               [Eq. (2)]
    3. P(t) = max over channels of pointiness traces
    4. tkeo(t) = max over channels of TKEO traces
    5. E_hpp = 0.6*zscore(P) + 0.4*zscore(tkeo)                          [Eq. (3)]
    6. E_hpp = gaussian_smooth(E_hpp, sigma=3)                            [Eq. (4)]

    # --- CET Evidence ---
    7. E_cet = HemiCET_UNet_ensemble(eeg)

    # --- Preprocessing ---
    8. Normalize E_hpp, E_cet to [0, 1]                                [Eqs. (5-6)]
    9. Threshold E_cet at max(80th percentile, 0.3)                    [Eqs. (7-8)]

    # --- Combination ---
   10. E(t) = max(E_hpp(t), E_cet(t)) + 3.0 * E_hpp(t) * E_cet(t)      [Eq. (9)]

    return E(t)
\end{verbatim}

\subsubsection{Output}

\begin{itemize}
  \item Combined evidence trace $E(t) \in \mathbb{R}^+$, used as input to the dynamic programming path optimizer for discharge time estimation.
\end{itemize}

\subsubsection{Performance}

The product-boosted max combination outperforms both individual evidence sources and simpler fusion strategies (e.g., weighted sum, max alone). The boost coefficient $\lambda = 3.0$ was selected via grid search over $\lambda \in \{0.5, 1.0, 2.0, 3.0, 4.0, 5.0\}$, optimizing discharge detection F1 score on a held-out validation set. The key advantage is that the method preserves the sensitivity of either individual detector (via the max term) while substantially boosting precision in regions of agreement (via the product term), yielding a favorable precision-recall tradeoff for downstream periodic discharge characterization.

\subsection{Dynamic programming with an approximately-periodic prior}
\label{app:dp}

\subsubsection{Overview}

We detect individual periodic discharges within each 10-second EEG segment using a Viterbi-style forward dynamic programming algorithm with an approximately-periodic prior. The algorithm receives a combined evidence trace $E(t)$ (from handcrafted and learned sources), identifies candidate discharge peaks, and selects the optimal approximately-periodic subsequence that balances evidence strength against temporal regularity. Post-processing includes EM-style template refinement and confidence-based filtering.

\subsubsection{Handcrafted evidence: HPP}

The handcrafted periodic pattern (HPP) evidence trace is computed per-channel from two complementary signal features.

\paragraph{Pointiness trace.}
For each local maximum at sample $n$ in the (optionally lowpass-filtered) signal $x$, we compute:
\[ \mathrm{prom}(n) = x(n) - \max\!\big(\min_{j \in [n-w, n)} x(j),\;\; \min_{j \in (n, n+w]} x(j)\big) \tag{1} \]
\[ \mathrm{width}(n) = \big|\{j : |j - n| \leq w,\; x(j) > x(n) - \tfrac{1}{2}\,\mathrm{prom}(n)\}\big| \tag{2} \]
\[ \mathrm{pt}(n) = \frac{\mathrm{prom}(n)^2}{\mathrm{width}(n)} \tag{3} \]
where $w = 8$ samples (40 ms) is the half-window for valley search, and $\mathrm{pt}(n) = 0$ at non-peak locations. Pointiness captures the sharpness of transient waveforms: narrow, prominent peaks (characteristic of discharges) receive high scores.

\paragraph{Teager-Kaiser energy operator (TKEO).}
The TKEO is applied to the 20 Hz lowpass-filtered signal $x_{\mathrm{lp}}$:
\[ \mathrm{tkeo}[n] = \big|x_{\mathrm{lp}}[n]^2 - x_{\mathrm{lp}}[n-1] \cdot x_{\mathrm{lp}}[n+1]\big| \tag{4} \]
TKEO responds to instantaneous energy and frequency, producing transient peaks at discharge locations.

\paragraph{Evidence combination.}
Both traces are z-scored within the 10-second window and combined with fixed weights:
\[ E_{\mathrm{hpp}}(t) = \max\!\Big(0,\;\; G_{\sigma} * \big[0.6 \cdot \tilde{\mathrm{pt}}(t) + 0.4 \cdot \widetilde{\mathrm{tkeo}}(t)\big]\Big) \tag{5} \]
where $\tilde{\cdot}$ denotes z-scoring, $G_{\sigma}$ is a Gaussian kernel with $\sigma = 3$ samples (15 ms), and $*$ denotes convolution. Negative values are clipped to zero.

\subsubsection{Neural evidence: CET-UNet}

Per-channel CET evidence $E_{\mathrm{cet}}(t)$ is obtained from the HemiCET-UNet ensemble (see HemiCET-UNet method).

\subsubsection{Evidence aggregation and combination}

\paragraph{Channel aggregation.}
Per-channel evidence is aggregated to a single trace depending on subtype:
\begin{itemize}
  \item \textbf{GPD:} median across all 18 channels
  \item \textbf{LPD with known laterality:} median across the 8 ipsilateral channels
  \item \textbf{LPD without laterality:} $\max\!\big(\mathrm{median}_{\mathcal{L}},\; \mathrm{median}_{\mathcal{R}}\big)$ pointwise
\end{itemize}

\paragraph{Product-boosted combination.}
The aggregated HPP and CET evidence traces are combined using a product-boost formula that amplifies regions where both sources agree:
\[ E(t) = \max\!\big(\hat{E}_{\mathrm{hpp}}(t),\; \hat{E}_{\mathrm{cet}}(t)\big) + \beta_{\mathrm{boost}} \cdot \hat{E}_{\mathrm{hpp}}(t) \cdot \hat{E}_{\mathrm{cet}}(t) \tag{6} \]
where $\hat{E}_{\mathrm{hpp}}$ and $\hat{E}_{\mathrm{cet}}$ are normalized to $[0, 1]$, the CET trace is thresholded at the 80th percentile of nonzero values (to suppress noise floor) and floored at 0.3, and $\beta_{\mathrm{boost}} = 3.0$.

\subsubsection{Active interval detection}

Before candidate extraction, we identify the longest contiguous interval of sustained discharge activity:
\begin{enumerate}
  \item Compute a rolling mean of $E(t)$ with a 1-second window ($w = 200$ samples)
  \item Threshold at 50\% of the global rolling maximum
  \item Find the longest contiguous above-threshold run
  \item Expand by 0.5 s on each side; require minimum 3 s duration (else use full segment)
\end{enumerate}

\subsubsection{Candidate extraction}

Candidate discharge peaks are extracted from $E(t)$ within the active interval via two passes:
\begin{enumerate}
  \item \textbf{Standard pass:} all local maxima with height $> 0.05 \cdot \max(E)$ and minimum inter-peak distance $0.2 \cdot T$ samples
  \item \textbf{Strong pass:} all local maxima with height $> 0.50 \cdot \max(E)$ and minimum distance $0.1 \cdot T$ samples
\end{enumerate}
Both sets are merged and deduplicated, yielding candidate set $\mathcal{S} = \{c_1, c_2, \ldots, c_n\}$ ordered by time.

\subsubsection{Dynamic programming formulation}

\paragraph{State space and scoring.}
Let $T = 1 / f_{\mathrm{est}}$ be the expected discharge period (from the CNN+ACF frequency ensemble). For each candidate $c \in \mathcal{S}$, define:

\textbf{Node score} (superlinear reward minus existence cost):
\[ R(c) = E(c)^{1.5} - \lambda \tag{7} \]
The exponent 1.5 provides superlinear scaling that disproportionately favors strong evidence peaks over marginal ones.

\textbf{Edge score} (transition from $c_i$ to $c_j$, allowing up to $M = 3$ skipped periods):
\[ Q(c_i, c_j) = \max_{m \in \{1, 2, 3\}} \left[-\alpha \cdot \left(\frac{\Delta t - m \cdot T}{m \cdot T}\right)^{\!2} - \beta \cdot (m - 1)\right] \tag{8} \]
where $\Delta t = (c_j - c_i) / F_s$ is the inter-candidate interval in seconds. The first term penalizes deviations from an integer multiple of the expected period (quadratic in fractional deviation). The second term penalizes skips: $m = 1$ incurs no penalty, $m = 2$ costs $\beta$, and $m = 3$ costs $2\beta$. The best skip count $m$ is chosen greedily.

\paragraph{Recurrence.}
The optimal score at candidate $c_j$ is:
\[ V(c_j) = \max\!\bigg(R(c_j),\;\; \max_{c_i < c_j,\; \Delta t \leq 4T} \big[V(c_i) + Q(c_i, c_j) + R(c_j)\big]\bigg) \tag{9} \]
with the first case allowing $c_j$ to start a new sequence.

\paragraph{Backtracking.}
The optimal path is recovered by backtracking from $c^* = \arg\max_j V(c_j)$.

\paragraph{Parameters.}
Parameter values are listed in table~\ref{tab:dp_params}.

\begin{table}[h]
\centering
\caption{Parameter values for the dynamic programming algorithm.}
\label{tab:dp_params}
\begin{tabular}{lll}
\hline
Parameter & Symbol & Value \\
\hline
Timing deviation penalty & $\alpha$ & 1.275 \\
Skip penalty & $\beta$ & 0.3 \\
Existence cost & $\lambda$ & 0.05 \\
Maximum skip & $M$ & 3 \\
\hline
\end{tabular}
\end{table}

\subsubsection{Algorithm pseudocode}

\begin{verbatim}
Algorithm: Forward DP with Approximately-Periodic Prior
Input: candidates S = {c_1, ..., c_n}, evidence E(t), freq_est f
Output: optimal discharge sequence D

  T = 1 / f                              // expected period
  For j = 1 to n:
      V[j] = E(c_j)^1.5 - lambda          // start new sequence
      prev[j] = -1

  For j = 2 to n:
      For i = 1 to j-1:
          dt = (c_j - c_i) / Fs
          if dt <= 0 or dt > 4T: continue
          edge = max over m in {1,2,3} of:
              -alpha * ((dt - m*T) / (m*T))^2 - beta * (m - 1)
          score = V[i] + edge + E(c_j)^1.5 - lambda
          if score > V[j]:
              V[j] = score
              prev[j] = i

  // Backtrack from best endpoint
  j* = argmax(V)
  D = []
  while j* >= 0:
      D.prepend(c_{j*})
      j* = prev[j*]
  return D
\end{verbatim}

\subsubsection{Post-processing}

\paragraph{EM template refinement.}
After the initial DP pass, discharge times are refined through template-based cross-correlation (3 iterations):
\begin{enumerate}
  \item \textbf{E-step:} Extract evidence snippets $E[d_k \pm 150\,\text{ms}]$ around each detected discharge $d_k$ and compute the mean template $\bar{\tau}$
  \item \textbf{M-step:} Compute normalized cross-correlation of $\bar{\tau}$ with the full evidence trace $E(t)$; extract new candidate peaks from the correlation signal
  \item Re-run the DP on the refined candidates
\end{enumerate}
\[ r(t) = \frac{\sum_{\delta} (\bar{\tau}(\delta) - \mu_{\bar{\tau}})(E(t + \delta) - \mu_{E,t})}{\|\bar{\tau} - \mu_{\bar{\tau}}\| \cdot \|E_t - \mu_{E,t}\|} \tag{10} \]

\paragraph{Post-hoc confidence filter.}
Detections with weak evidence are removed:
\[ \mathcal{D}_{\mathrm{filtered}} = \big\{d \in \mathcal{D} : E(d) \geq \eta \cdot \mathrm{median}_{d' \in \mathcal{D}} E(d')\big\} \tag{11} \]
with $\eta = 0.3$ (minimum 30\% of the median peak evidence).

\subsubsection{Final frequency estimate}

The output frequency is derived from the inter-peak intervals (IPIs) of the final discharge sequence:
\[ f_{\mathrm{IPI}} = \frac{1}{\mathrm{median}\!\big(\{d_{k+1} - d_k\}_{k=1}^{K-1}\big)} \tag{12} \]
where $d_k$ are discharge times in seconds.

\paragraph{Per-channel timing.}
Global discharge times are refined per-channel by searching for the local evidence maximum within $\pm 50$ ms of each global time, accommodating inter-channel propagation delays.

\subsubsection{Performance}

The full pipeline (product-boosted HPP+CET evidence, DP with CNN+ACF frequency prior, EM refinement, post-hoc filtering) achieves discharge detection F1 of approximately 0.74 at $\pm 100$ ms tolerance, evaluated on 593 patients with expert-reviewed discharge annotations. The IPI-derived frequency achieves Spearman $\rho \approx 0.65$ with expert frequency ratings.

\subsection{CNN+ACF frequency ensemble}
\label{app:cnn_acf}

\subsubsection{Overview}

Accurate estimation of the discharge repetition frequency is critical both as a clinical descriptor and as the periodic prior $T = 1/f$ that governs the dynamic programming discharge detector. We combine two complementary frequency estimates --- one learned (CNN) and one handcrafted (autocorrelation of the pointiness trace) --- into a robust ensemble that serves as the frequency prior for downstream processing.

\subsubsection{CNN frequency estimate}

\paragraph{Per-channel prediction.}
Each of the 18 bipolar channels is processed independently through the 5-fold ChannelPD-Net ensemble (see ChannelPD-Net method), yielding per-channel PD probability $p_c$ and log-frequency prediction $\hat{f}_{\log,c}$.

\paragraph{PD-weighted aggregation.}
The patient-level CNN frequency is a PD-weighted average in log-space, which down-weights channels lacking periodic patterns:
\[ f_{\mathrm{CNN}} = \exp\!\left(\frac{\sum_{c=1}^{C} p_c \cdot \hat{f}_{\log,c}}{\sum_{c=1}^{C} p_c}\right) \tag{1} \]
where $C = 18$ is the number of bipolar channels. If $\sum_c p_c < 10^{-6}$, the unweighted mean is used as a fallback. The estimate is clipped to the physiological range $[0.3, 3.5]$ Hz.

\textbf{Rationale for log-space averaging.} Discharge frequencies span nearly an order of magnitude (0.3--3.5 Hz). Averaging in log-space prevents high-frequency outliers from dominating and provides a more symmetric error distribution.

\subsubsection{ACF frequency estimate}

\paragraph{Per-channel autocorrelation.}
For each channel, we compute the autocorrelation of the pointiness trace $\mathrm{pt}(t)$ (see Dynamic Programming method, Eq.~3) via the Wiener-Khinchin theorem:
\[ R(\tau) = \mathcal{F}^{-1}\!\Big\{\big|\mathcal{F}\{\mathrm{pt}(t) - \bar{\mathrm{pt}}\}\big|^2\Big\} \tag{2} \]
The autocorrelation is normalized by its zero-lag value:
\[ \hat{R}(\tau) = \frac{R(\tau)}{R(0)} \tag{3} \]
so that $\hat{R}(0) = 1$.

\paragraph{Peak detection.}
The first prominent peak in the normalized ACF within the physiologically plausible lag range indicates the dominant repetition period:
\[ \tau^* = \arg\max_{\tau \in [\tau_{\min},\, \tau_{\max}]} \hat{R}(\tau), \quad \text{subject to } \hat{R}(\tau^*) > 0.1 \tag{4} \]
where $\tau_{\min} = \lfloor F_s / 3.5 \rfloor$ samples (maximum 3.5 Hz) and $\tau_{\max} = \lfloor F_s / 0.33 \rfloor$ samples (minimum 0.33 Hz). If multiple peaks exist, the shortest-lag peak above threshold is selected to avoid subharmonic aliasing.

The per-channel ACF frequency is:
\[ f_{\mathrm{ACF}}^{(c)} = \frac{F_s}{\tau^*_c} \tag{5} \]

\paragraph{Multi-channel aggregation.}
The patient-level ACF frequency is the median of valid per-channel estimates:
\[ f_{\mathrm{ACF}} = \mathrm{median}\!\big(\{f_{\mathrm{ACF}}^{(c)} : f_{\mathrm{ACF}}^{(c)} \text{ is finite}\}\big) \tag{6} \]
The median provides robustness against channels with artifact or absent periodic patterns. If no channel yields a valid ACF peak, $f_{\mathrm{ACF}}$ is treated as missing.

\subsubsection{Ensemble combination}

The final frequency prior combines the two estimates with fixed weights favoring the CNN:
\[ f_{\mathrm{prior}} = \begin{cases} 0.8 \cdot f_{\mathrm{CNN}} + 0.2 \cdot f_{\mathrm{ACF}} & \text{if } f_{\mathrm{ACF}} \text{ is finite} \\ f_{\mathrm{CNN}} & \text{otherwise} \end{cases} \tag{7} \]
The combined estimate is clipped to the physiological range:
\[ f_{\mathrm{prior}} = \mathrm{clip}(f_{\mathrm{prior}},\; 0.3,\; 3.5) \tag{8} \]
This $f_{\mathrm{prior}}$ determines the expected period $T = 1 / f_{\mathrm{prior}}$ used as the periodic prior in the DP algorithm (see Dynamic Programming method, Eq.~8).

\subsubsection{Algorithm pseudocode}

\begin{verbatim}
Algorithm: CNN+ACF Frequency Ensemble
Input: segment_18ch (18 x 2000), subtype, laterality
Output: f_prior (Hz)

  // CNN frequency
  For each channel c = 1, ..., 18:
      z-score normalize channel c
      For each fold k = 1, ..., 5:
          (p_c^k, f_log_c^k, _) = ChannelPDNet_k(channel_c)
      p_c = mean(p_c^1, ..., p_c^5)
      f_log_c = mean(f_log_c^1, ..., f_log_c^5)
  f_CNN = exp( sum(p_c * f_log_c) / sum(p_c) )
  f_CNN = clip(f_CNN, 0.3, 3.5)

  // ACF frequency
  acf_freqs = []
  For each channel c = 1, ..., 18:
      pt_c = pointiness_trace(lowpass_20Hz(channel_c))
      R_c = IFFT(|FFT(pt_c - mean(pt_c))|^2)
      R_c = R_c / R_c[0]
      Find first peak tau* in R_c[Fs/3.5 : Fs/0.33] with height > 0.1
      if peak found:
          acf_freqs.append(Fs / tau*)
  f_ACF = median(acf_freqs) if len(acf_freqs) > 0 else NaN

  // Ensemble
  if f_ACF is finite:
      f_prior = 0.8 * f_CNN + 0.2 * f_ACF
  else:
      f_prior = f_CNN
  f_prior = clip(f_prior, 0.3, 3.5)
  return f_prior
\end{verbatim}

\subsubsection{Design rationale}

The 80/20 weighting reflects empirical findings that the CNN frequency estimate is more accurate on average (having been trained on expert labels), while the ACF provides a useful corrective signal in cases where the CNN mispredicts --- particularly for segments with unusual morphology or low PD probability across channels. The ACF is parameter-free and does not require labeled data, making it a complementary regularizer.

\subsubsection{Performance}

\begin{table}[h]
\centering
\caption{Spearman correlation of CNN, ACF, and ensemble frequency estimates with expert ratings.}
\label{tab:cnn_acf_perf}
\begin{tabular}{ll}
\hline
Method & Spearman $\rho$ vs.\ expert \\
\hline
CNN alone & $\sim$0.55 \\
ACF alone & $\sim$0.40 \\
CNN+ACF ensemble (0.8/0.2) & $\sim$0.60 \\
Full pipeline (IPI from DP) & $\sim$0.65 \\
\hline
\end{tabular}
\end{table}

The CNN+ACF ensemble provides the prior that seeds the DP algorithm; the final IPI-derived frequency from the DP output further improves correlation with expert ratings by leveraging the detected discharge times directly.

\subsection{Discharge-locked topographic localization}
\label{app:discharge_topo}

\subsubsection{Overview}

We compute mean discharge topographies by aligning monopolar EEG epochs to discharge peaks using a two-pass Laplacian-GFP refinement procedure. The first pass estimates a mean template from coarsely aligned epochs; the second pass cross-correlates individual epoch GFP profiles against the template to correct residual jitter. A GFP-squared weighting scheme suppresses phantom discharges---timepoints where dynamic programming interpolated a discharge but no true electrographic transient exists. The resulting topography is rendered as a spherical spline interpolation and verbally described via a 16-region parcellation.

\subsubsection{Input}

\begin{itemize}
  \item 19-channel monopolar EEG, $V[e,t]$, bandpass filtered 0.5--20 Hz, where $e \in \{1,\ldots,19\}$ indexes electrodes in the International 10--20 system and $t$ is the sample index.
  \item Discharge times $\{t_1, t_2, \ldots, t_K\}$ from the HemiCET+DP pipeline.
\end{itemize}

\subsubsection{Mathematical formulation}

\paragraph{Step 1: Surface Laplacian.}
The surface Laplacian sharpens spatial resolution by removing volume-conducted contributions. For each electrode $e$ with neighbor set $\mathcal{N}_e$ in the 10--20 montage adjacency graph:
\[ V_{\mathrm{lap}}[e,t] = V[e,t] - \frac{1}{|\mathcal{N}_e|} \sum_{n \in \mathcal{N}_e} V[n,t] \tag{1} \]

\paragraph{Step 2: Global field power.}
The global field power (GFP) quantifies the instantaneous spatial standard deviation across all $C = 19$ channels:
\[ \mathrm{gfp}(t) = \sqrt{\frac{1}{C} \sum_{e=1}^{C} V_{\mathrm{lap}}[e,t]^2} \tag{2} \]
GFP peaks correspond to moments of maximal scalp field strength and provide a physiologically motivated alignment target for discharge epochs.

\paragraph{Step 3: Pass 1 --- coarse alignment and template extraction.}
For each discharge time $t_k$, locate the GFP peak within a $\pm 25\,\mathrm{ms}$ window:
\[ t_k^* = \underset{|t - t_k| \leq 25\,\mathrm{ms}}{\operatorname{argmax}}\; \mathrm{gfp}(t) \tag{3} \]
Extract a $\pm 50\,\mathrm{ms}$ epoch centered at $t_k^*$:
\[ \mathbf{E}_k[\tau] = V_{\mathrm{lap}}[\cdot,\; t_k^* + \tau], \quad \tau \in [-50\,\mathrm{ms},\; +50\,\mathrm{ms}] \tag{4} \]
The template GFP profile is the epoch-averaged GFP:
\[ g_{\mathrm{tmpl}}(\tau) = \frac{1}{K} \sum_{k=1}^{K} \mathrm{gfp}(t_k^* + \tau) \tag{5} \]

\paragraph{Step 4: Pass 2 --- cross-correlation refinement.}
For each epoch $k$, compute the normalized cross-correlation between its GFP profile $g_k(\tau) = \mathrm{gfp}(t_k^* + \tau)$ and the template $g_{\mathrm{tmpl}}(\tau)$:
\[ R_k(\delta) = \frac{\sum_{\tau} g_k(\tau + \delta)\, g_{\mathrm{tmpl}}(\tau)}{\sqrt{\sum_{\tau} g_k(\tau+\delta)^2 \;\cdot\; \sum_{\tau} g_{\mathrm{tmpl}}(\tau)^2}} \tag{6} \]
The refined alignment shift is:
\[ \delta_k^* = \underset{|\delta| \leq 25\,\mathrm{ms}}{\operatorname{argmax}}\; R_k(\delta) \tag{7} \]
The refined peak time becomes $\hat{t}_k = t_k^* + \delta_k^*$, and the voltage vector at the refined peak is $\mathbf{v}_k = V_{\mathrm{lap}}[\cdot,\; \hat{t}_k] \in \mathbb{R}^{19}$.

\paragraph{Step 5: GFP-squared weighted averaging.}
The mean topography is computed as a weighted average with $w_k = \mathrm{gfp}(\hat{t}_k)^2$:
\[ \mathrm{topo}[e] = \frac{\sum_{k=1}^{K} w_k \cdot V_{\mathrm{lap}}[e,\; \hat{t}_k]}{\sum_{k=1}^{K} w_k} \tag{8} \]
The squared GFP weighting ensures that epochs with strong, well-defined spatial patterns dominate the average, while phantom discharges---which arise from dynamic programming interpolation at timepoints lacking genuine electrographic transients---contribute minimally due to their low GFP.

\subsubsection{Algorithm}

\begin{verbatim}
DISCHARGE_TOPOGRAPHY(V[e,t], {t_1,...,t_K}):
    1. Compute Laplacian: V_lap[e,t] via Eq. (1)
    2. Compute GFP: gfp(t) via Eq. (2)
    3. PASS 1 --- Coarse alignment:
       for k = 1 to K:
           t*_k = argmax gfp(t) for |t - t_k| <= 25ms     [Eq. (3)]
           Extract epoch E_k[tau] for tau in [-50ms, +50ms]  [Eq. (4)]
       Compute template GFP profile g_tmpl(tau)              [Eq. (5)]
    4. PASS 2 --- Refined alignment:
       for k = 1 to K:
           Compute cross-correlation R_k(delta)              [Eq. (6)]
           delta*_k = argmax R_k(delta) for |delta| <= 25ms  [Eq. (7)]
           t_hat_k = t*_k + delta*_k
           v_k = V_lap[:, t_hat_k]
           w_k = gfp(t_hat_k)^2
    5. Compute weighted mean topography via Eq. (8)
    6. Render via MNE spherical spline topoplot (inferno colormap)
    7. Generate verbal description via describe_ied_topoplot()
    return topo[e]
\end{verbatim}

\subsubsection{Rendering and description}

The 19-electrode topography vector $\mathrm{topo}[e]$ is rendered using MNE-Python's spherical spline interpolation onto a standard head model with the \emph{inferno} colormap (perceptually uniform, colorblind-safe). Verbal spatial descriptions are generated by the \texttt{describe\_ied\_topoplot()} function in \texttt{morgoth-viewer}, which maps electrode amplitudes to 16 anatomical brain regions (e.g., left frontal, right temporal, midline parietal) and reports the top regions by normalized amplitude.

\subsubsection{Performance}

The two-pass alignment reduces cross-epoch jitter by approximately 8--12 ms on average compared to single-pass GFP alignment, yielding sharper mean topographies with higher peak GFP. The GFP-squared weighting reduces the effective contribution of phantom discharges by a factor of 4--10 relative to uniform averaging.

\subsection{Hybrid-PLV PD spatial extent}
\label{app:hybrid_plv}

\subsubsection{Overview}

Spatial localization of periodic discharges requires identifying which brain regions participate in the discharge pattern. We employ a hybrid approach that combines per-channel PD probabilities from ChannelPD-Net with phase-locking value (PLV) analysis to produce region-level involvement scores. The CNN selects reference channels most likely to contain discharges, and PLV quantifies phase coherence between each channel and this reference, identifying regions that are synchronously involved in the periodic pattern.

\subsubsection{Reference channel selection}

The top-3 channels by ChannelPD-Net probability are selected as references. For LPD segments with known laterality, contralateral channels are suppressed before ranking to ensure the reference is seeded from the correct hemisphere:
\[ \tilde{p}_c = \begin{cases} p_c & \text{if } c \in \mathcal{H}_{\mathrm{ipsi}} \text{ or subtype is GPD} \\ 0 & \text{if } c \in \mathcal{H}_{\mathrm{contra}} \text{ and subtype is LPD} \end{cases} \tag{1} \]
\[ \mathcal{R} = \operatorname{argtop\text{-}3}(\tilde{p}_c) \tag{2} \]
where $\mathcal{H}_{\mathrm{ipsi}}$ and $\mathcal{H}_{\mathrm{contra}}$ are the ipsilateral and contralateral hemisphere channel index sets, and $\mathcal{R}$ is the set of three reference channels.

\subsubsection{Bandpass filtering}

All 18 channels are bandpass-filtered to the periodic discharge frequency range using a 4th-order Butterworth filter:
\[ x_c^{\mathrm{bp}}(t) = \mathrm{BPF}_{[0.5,\, 3.5]\,\text{Hz}}\!\big(x_c(t)\big) \tag{3} \]
This isolates the slow periodic component while rejecting faster EEG rhythms and DC drift.

\subsubsection{Phase extraction via Hilbert transform}

The instantaneous phase of each channel is obtained from the analytic signal:
\[ z_c(t) = x_c^{\mathrm{bp}}(t) + j \cdot \mathcal{H}\!\big\{x_c^{\mathrm{bp}}(t)\big\} \tag{4} \]
\[ \varphi_c(t) = \arg\!\big(z_c(t)\big) \in (-\pi, \pi] \tag{5} \]
where $\mathcal{H}\{\cdot\}$ denotes the Hilbert transform.

\subsubsection{Reference phase construction}

The reference phase is constructed as the circular mean of the three reference channel phases:
\[ \varphi_{\mathrm{ref}}(t) = \arg\!\left(\sum_{k \in \mathcal{R}} \exp\!\big(j \cdot \varphi_k(t)\big)\right) \tag{6} \]
This produces a robust reference that is less sensitive to noise in any single channel than using a single reference.

\subsubsection{Per-channel phase-locking value}

The PLV between each channel and the reference quantifies the consistency of their phase relationship across time:
\[ \mathrm{PLV}_c = \left|\frac{1}{N} \sum_{t=1}^{N} \exp\!\Big(j \cdot \big(\varphi_c(t) - \varphi_{\mathrm{ref}}(t)\big)\Big)\right| \in [0, 1] \tag{7} \]
where $N = 2000$ is the number of samples. $\mathrm{PLV}_c = 1$ indicates perfect phase synchrony (constant phase difference), while $\mathrm{PLV}_c \approx 0$ indicates no consistent phase relationship.

\subsubsection{Combined channel score}

The CNN probability and PLV are combined with equal weight to produce a per-channel involvement score:
\[ s_c = 0.5 \cdot p_c + 0.5 \cdot \mathrm{PLV}_c \tag{8} \]
This combination leverages the CNN's ability to detect discharge morphology (even without phase information) and the PLV's sensitivity to temporal synchrony (even in channels with atypical morphology).

\subsubsection{Region mapping}

The 18 bipolar channels are mapped to 8 anatomical regions corresponding to the standard double-banana montage; see table~\ref{tab:hybrid_plv_regions}.

\begin{table}[h]
\centering
\caption{Mapping of bipolar channels to anatomical regions for the Hybrid-PLV spatial extent computation.}
\label{tab:hybrid_plv_regions}
\begin{tabular}{lll}
\hline
Region & Abbreviation & Channel indices \\
\hline
Left Frontal & LF & Fp1-F7 (0), Fp1-F3 (8) \\
Right Frontal & RF & Fp2-F8 (4), Fp2-F4 (12) \\
Left Temporal & LT & F7-T3 (1), T3-T5 (2) \\
Right Temporal & RT & F8-T4 (5), T4-T6 (6) \\
Left Centro-Parietal & LCP & F3-C3 (9), C3-P3 (10) \\
Right Centro-Parietal & RCP & F4-C4 (13), C4-P4 (14) \\
Left Occipital & LO & T5-O1 (3), P3-O1 (11) \\
Right Occipital & RO & T6-O2 (7), P4-O2 (15) \\
\hline
\end{tabular}
\end{table}

The region score is the maximum channel score among channels belonging to that region:
\[ S_r = \max_{c \in \mathcal{C}_r} s_c, \quad r \in \{\text{LF, RF, LT, RT, LCP, RCP, LO, RO}\} \tag{9} \]
where $\mathcal{C}_r$ is the set of channel indices mapped to region $r$. Taking the maximum (rather than mean) ensures that a region is considered involved if any of its constituent channels shows strong evidence.

\subsubsection{Binary involvement and spatial extent}

Each region is classified as involved if its score exceeds a fixed threshold:
\[ \mathrm{involved}(r) = \mathbb{1}[S_r > \theta], \quad \theta = 0.4 \tag{10} \]
The spatial extent of the discharge pattern is the fraction of regions involved:
\[ \mathrm{extent} = \frac{\sum_{r=1}^{8} \mathbb{1}[S_r > \theta]}{8} \in [0, 1] \tag{11} \]

\subsubsection{Algorithm pseudocode}

\begin{verbatim}
Algorithm: Hybrid-PLV Spatial Localization
Input: segment_18ch (18 x 2000), channel_probs p (18,), laterality
Output: region_scores (8,), involved_regions, spatial_extent

  // 1. Reference selection
  if LPD and laterality known:
      suppress contralateral channel probs to 0
  ref_channels = argsort(p, descending)[:3]

  // 2. Bandpass filter 0.5-3.5 Hz
  For each channel c = 1, ..., 18:
      x_bp[c] = butterworth_bandpass(segment_18ch[c], 0.5, 3.5 Hz)

  // 3. Phase extraction
  For each channel c = 1, ..., 18:
      phi[c] = angle(hilbert(x_bp[c]))

  // 4. Reference phase (circular mean of top-3)
  phi_ref = angle( sum_{k in ref_channels} exp(j * phi[k]) )

  // 5. PLV computation
  For each channel c = 1, ..., 18:
      PLV[c] = |mean_t( exp(j * (phi[c] - phi_ref)) )|
      s[c] = 0.5 * p[c] + 0.5 * PLV[c]

  // 6. Region mapping
  For each region r in {LF, RF, LT, RT, LCP, RCP, LO, RO}:
      S[r] = max(s[c] for c in region_channels[r])
      involved[r] = (S[r] > 0.4)

  // 7. Spatial extent
  extent = sum(involved) / 8

  return S, involved, extent
\end{verbatim}

\subsubsection{Design rationale}

Phase-locking value was selected from a 26-method spatial localization contest (including amplitude correlation, coherence, Granger causality, CSD analysis, and others) as the best-performing single spatial metric. The hybrid combination with CNN probabilities outperforms either component alone: the CNN provides morphology-based evidence that is robust to phase noise, while PLV captures temporal synchrony patterns that the CNN's channel-independent architecture cannot directly model.

The 0.5/0.5 weighting was determined empirically. The threshold $\theta = 0.4$ was chosen to balance sensitivity and specificity for region involvement classification.

\subsubsection{Performance}

The Hybrid-PLV method achieves a composite spatial localization score of approximately 0.79 and a region-level AUC of approximately 0.71, evaluated against expert spatial annotations. PLV alone achieves AUC approximately 0.68, while CNN probabilities alone achieve AUC approximately 0.65, confirming the complementary value of the two information sources.

\subsection{NB-Hilbert: iterative narrowband Hilbert refinement for RDA}
\label{app:nb_hilbert}

\subsubsection{Overview}

We estimate the frequency and lateralization of rhythmic delta activity (LRDA/GRDA) using a two-pass iterative narrowband refinement procedure. The first pass applies a broad delta bandpass to obtain coarse estimates of dominant hemisphere and frequency via Hilbert-based instantaneous frequency analysis. The second pass re-filters the signal in a narrow band centered on the first-pass frequency estimate, yielding refined lateralization scores and a more precise frequency measurement. The method also classifies each segment as LRDA or GRDA based on an amplitude asymmetry index.

\subsubsection{Input}

\begin{itemize}
  \item 18-channel bipolar EEG, $\mathbf{x}(t) \in \mathbb{R}^{18}$, sampled at $f_s = 200\,\mathrm{Hz}$, with segment duration $T = 10\,\mathrm{s}$.
  \item Channels are partitioned into left-hemisphere set $\mathcal{L}$ and right-hemisphere set $\mathcal{R}$.
\end{itemize}

\subsubsection{Mathematical formulation}

\paragraph{Pass 1: coarse estimation --- broadband delta filtering.}
Apply a bandpass filter at 0.5--3.5 Hz (3rd-order Butterworth, zero-phase) to all channels:
\[ x_{\mathrm{broad}}^{(c)}(t) = \mathrm{BPF}_{[0.5,\, 3.5]}\!\left[x^{(c)}(t)\right], \quad c \in \{1,\ldots,18\} \tag{1} \]

\paragraph{Coarse lateralization.}
Compute the mean variance for each hemisphere:
\[ \sigma_L^2 = \frac{1}{|\mathcal{L}|} \sum_{c \in \mathcal{L}} \mathrm{var}\!\left(x_{\mathrm{broad}}^{(c)}\right), \qquad \sigma_R^2 = \frac{1}{|\mathcal{R}|} \sum_{c \in \mathcal{R}} \mathrm{var}\!\left(x_{\mathrm{broad}}^{(c)}\right) \tag{2} \]
The dominant hemisphere is:
\[ \mathrm{dom} = \underset{h \in \{L, R\}}{\operatorname{argmax}}\; \sigma_h^2 \tag{3} \]
Select the top-3 channels by variance from the dominant hemisphere: $\{c_1, c_2, c_3\} \subset \mathcal{H}_{\mathrm{dom}}$.

\paragraph{Hilbert-based instantaneous frequency.}
For each selected channel $c_i$, compute the analytic signal:
\[ z^{(c_i)}(t) = x_{\mathrm{broad}}^{(c_i)}(t) + j \cdot \mathcal{H}\!\left\{x_{\mathrm{broad}}^{(c_i)}\right\}(t) \tag{4} \]
where $\mathcal{H}\{\cdot\}$ denotes the Hilbert transform. The instantaneous phase and frequency are:
\[ \phi^{(c_i)}(t) = \arg\!\left(z^{(c_i)}(t)\right) \tag{5} \]
\[ f^{(c_i)}(t) = \frac{f_s}{2\pi} \cdot \Delta\!\left(\mathrm{unwrap}\!\left(\phi^{(c_i)}\right)\right)(t) \tag{6} \]
where $\Delta$ denotes the first-order backward difference operator. The coarse frequency estimate is obtained by pooling valid instantaneous frequency samples across the three channels:
\[ \hat{f}_{\mathrm{coarse}} = \underset{t,\, i}{\mathrm{median}}\;\left\{f^{(c_i)}(t) : f^{(c_i)}(t) \in [0.3,\; 4.0]\,\mathrm{Hz}\right\} \tag{7} \]

\paragraph{Pass 2: narrowband refinement --- narrowband filtering.}
Re-filter all channels using a narrow bandpass centered on the coarse estimate:
\[ x_{\mathrm{narrow}}^{(c)}(t) = \mathrm{BPF}_{[\hat{f}_{\mathrm{coarse}} - 0.4,\;\; \hat{f}_{\mathrm{coarse}} + 0.4]}\!\left[x^{(c)}(t)\right] \tag{8} \]

\paragraph{Refined lateralization.}
Compute the analytic signal for each narrowband channel and extract the instantaneous envelope:
\[ \mathrm{env}^{(c)}(t) = \left|z_{\mathrm{narrow}}^{(c)}(t)\right| = \left|x_{\mathrm{narrow}}^{(c)}(t) + j \cdot \mathcal{H}\!\left\{x_{\mathrm{narrow}}^{(c)}\right\}(t)\right| \tag{9} \]
Hemisphere amplitude scores:
\[ L_{\mathrm{score}} = \frac{1}{|\mathcal{L}|} \sum_{c \in \mathcal{L}} \overline{\mathrm{env}^{(c)}}, \qquad R_{\mathrm{score}} = \frac{1}{|\mathcal{R}|} \sum_{c \in \mathcal{R}} \overline{\mathrm{env}^{(c)}} \tag{10} \]
where $\overline{\mathrm{env}^{(c)}} = \frac{1}{T} \int_0^T \mathrm{env}^{(c)}(t)\, dt$ is the temporal mean envelope amplitude.

The refined dominant hemisphere is $\mathrm{dom}' = \operatorname{argmax}_{h}(L_{\mathrm{score}}, R_{\mathrm{score}})$.

\paragraph{Refined frequency.}
Select the top-3 narrowband channels by variance on the refined dominant hemisphere. Repeat the Hilbert instantaneous frequency estimation (Eqs.~4--7) on these narrowband signals to obtain:
\[ \hat{f}_{\mathrm{refined}} = \underset{t,\, i}{\mathrm{median}}\;\left\{f_{\mathrm{narrow}}^{(c_i)}(t) : f_{\mathrm{narrow}}^{(c_i)}(t) \in [0.3,\; 4.0]\,\mathrm{Hz}\right\} \tag{11} \]

\paragraph{LRDA vs.\ GRDA classification.}
Compute the asymmetry index:
\[ A = \frac{|L_{\mathrm{score}} - R_{\mathrm{score}}|}{L_{\mathrm{score}} + R_{\mathrm{score}}} \tag{12} \]
The subtype classification is:
\[ \mathrm{subtype} = \begin{cases} \text{LRDA} & \text{if } A > \theta_A \\ \text{GRDA} & \text{otherwise} \end{cases} \tag{13} \]
where $\theta_A$ is a threshold calibrated against expert consensus labels.

\subsubsection{Algorithm}

\begin{verbatim}
W05_RDA(x[c,t], f_s=200):
    # --- PASS 1: Coarse ---
    1. x_broad[c,t] = BPF(x[c,t], [0.5, 3.5] Hz)           [Eq. (1)]
    2. Compute sigma_L^2, sigma_R^2                           [Eq. (2)]
    3. dom = argmax(sigma_L^2, sigma_R^2)                     [Eq. (3)]
    4. Select top-3 channels {c1,c2,c3} by variance on dom
    5. For each c_i:
         z(t) = x_broad(t) + j*Hilbert(x_broad(t))           [Eq. (4)]
         phi(t) = arg(z(t))                                   [Eq. (5)]
         f(t) = (f_s/2pi) * diff(unwrap(phi))                 [Eq. (6)]
    6. f_coarse = median({f(t) : f(t) in [0.3, 4.0]})        [Eq. (7)]

    # --- PASS 2: Narrowband ---
    7. x_narrow[c,t] = BPF(x[c,t], [f_coarse-0.4, f_coarse+0.4])  [Eq. (8)]
    8. env[c](t) = |hilbert(x_narrow[c,t])|                   [Eq. (9)]
    9. L_score = mean(env_left), R_score = mean(env_right)     [Eq. (10)]
   10. dom' = argmax(L_score, R_score)
   11. Select top-3 narrowband channels on dom'
   12. f_refined = Hilbert freq on narrowband top-3             [Eq. (11)]
   13. A = |L - R| / (L + R)                                   [Eq. (12)]
   14. subtype = LRDA if A > theta_A else GRDA                 [Eq. (13)]

    return subtype, laterality(dom'), f_refined
\end{verbatim}

\subsubsection{Output}

\begin{itemize}
  \item \textbf{Subtype}: LRDA or GRDA
  \item \textbf{Laterality}: left or right dominant hemisphere (for LRDA)
  \item \textbf{Frequency}: refined estimate $\hat{f}_{\mathrm{refined}}$ in Hz
\end{itemize}

\subsubsection{Performance}

The two-pass approach improves frequency estimation accuracy by 15--20\% relative to single-pass broadband Hilbert analysis, as measured by Spearman correlation against expert-labeled frequencies. The narrowband filtering in Pass 2 suppresses harmonic contamination and adjacent-band activity that bias instantaneous frequency estimates in the broadband case. In the V5 lateralization contest (76 methods), W05 achieved the best unified weighted AUC of 0.837 with frequency Spearman $\rho = 0.635$.

\subsection{RDA-PLV spatial extent}
\label{app:rda_plv}

\subsubsection{Overview}

We quantify the spatial extent of rhythmic delta activity (RDA) using a phase-locking value (PLV) metric weighted by narrowband amplitude. Each bipolar channel is compared to a reference signal derived from the dominant-hemisphere channels in the narrowband around the estimated RDA frequency. Channels whose PLV-amplitude product exceeds a threshold are classified as ``involved,'' and spatial extent is defined as the fraction of involved channels. This approach captures both phase coherence (indicating true rhythmic coupling) and amplitude (indicating signal strength), avoiding false positives from low-amplitude channels with spuriously high PLV.

\subsubsection{Input}

\begin{itemize}
  \item 18-channel bipolar EEG, $\mathbf{x}(t) \in \mathbb{R}^{18}$
  \item Estimated RDA frequency $\hat{f}\,(\mathrm{Hz})$ from the W05 pipeline (Method 7)
  \item Dominant hemisphere and top-3 channels from W05
\end{itemize}

\subsubsection{Mathematical formulation}

\paragraph{Step 1: narrowband filtering.}
Apply a 3rd-order Butterworth bandpass filter centered on the estimated frequency:
\[ x_{\mathrm{nb}}^{(c)}(t) = \mathrm{BPF}_{[\hat{f} - 0.4,\;\; \hat{f} + 0.4]}^{(\mathrm{Butter},\, 3)}\!\left[x^{(c)}(t)\right], \quad c \in \{1,\ldots,18\} \tag{1} \]

\paragraph{Step 2: reference signal construction.}
The reference signal is the arithmetic mean of the narrowband-filtered top-3 channels on the dominant hemisphere, $\{c_1, c_2, c_3\}$:
\[ x_{\mathrm{ref}}(t) = \frac{1}{3} \sum_{i=1}^{3} x_{\mathrm{nb}}^{(c_i)}(t) \tag{2} \]

\paragraph{Step 3: phase extraction.}
Compute the instantaneous phase for each channel and the reference via the Hilbert transform:
\[ \phi^{(c)}(t) = \arg\!\left(x_{\mathrm{nb}}^{(c)}(t) + j \cdot \mathcal{H}\!\left\{x_{\mathrm{nb}}^{(c)}\right\}(t)\right) \tag{3} \]
\[ \phi_{\mathrm{ref}}(t) = \arg\!\left(x_{\mathrm{ref}}(t) + j \cdot \mathcal{H}\!\left\{x_{\mathrm{ref}}\right\}(t)\right) \tag{4} \]

\paragraph{Step 4: phase-locking value.}
The PLV for each channel quantifies the consistency of the phase difference relative to the reference over $N$ time samples:
\[ \mathrm{PLV}_c = \left|\frac{1}{N} \sum_{t=1}^{N} \exp\!\left(j \cdot \left[\phi^{(c)}(t) - \phi_{\mathrm{ref}}(t)\right]\right)\right| \tag{5} \]
$\mathrm{PLV}_c \in [0,1]$, where 1 indicates perfect phase locking and 0 indicates uniformly distributed phase differences.

\paragraph{Step 5: amplitude ratio.}
Compute the mean envelope amplitude per channel and normalize by the maximum across channels:
\[ \bar{a}_c = \frac{1}{N} \sum_{t=1}^{N} \left|x_{\mathrm{nb}}^{(c)}(t) + j \cdot \mathcal{H}\!\left\{x_{\mathrm{nb}}^{(c)}\right\}(t)\right| \tag{6} \]
\[ a_c = \frac{\bar{a}_c}{\max_{c'} \bar{a}_{c'}} \tag{7} \]

\paragraph{Step 6: combined score.}
The per-channel involvement score is the product of PLV and normalized amplitude:
\[ s_c = \mathrm{PLV}_c \cdot a_c \tag{8} \]
This product penalizes channels that have high PLV but negligible amplitude (noise-driven phase locking) and channels with high amplitude but poor phase coherence (non-rhythmic activity).

\paragraph{Step 7: variance explained.}
As an auxiliary quality metric, fit a sinusoidal model to each narrowband channel:
\[ \hat{x}^{(c)}(t) = A_c \sin(2\pi \hat{f} \cdot t + \phi_{0,c}) + B_c \tag{9} \]
where $A_c$, $\phi_{0,c}$, and $B_c$ are estimated via least squares. The variance explained is:
\[ \mathrm{VE}_c = 1 - \frac{\mathrm{var}\!\left(x_{\mathrm{nb}}^{(c)} - \hat{x}^{(c)}\right)}{\mathrm{var}\!\left(x_{\mathrm{nb}}^{(c)}\right)} \tag{10} \]

\paragraph{Step 8: thresholding and spatial extent.}
A channel is classified as involved if its combined score exceeds a fixed threshold:
\[ \mathcal{I} = \left\{c : s_c > \theta_s\right\}, \quad \theta_s = 0.62 \tag{11} \]
The spatial extent is:
\[ \mathrm{SE} = \frac{|\mathcal{I}|}{18} \tag{12} \]
where the denominator is the total number of bipolar channels.

\subsubsection{Algorithm}

\begin{verbatim}
RDA_PLV(x[c,t], f_hat, top3_channels):
    1. x_nb[c,t] = BPF(x[c,t], [f_hat-0.4, f_hat+0.4], order=3)   [Eq. (1)]
    2. x_ref(t) = mean(x_nb[top3_channels, t])                      [Eq. (2)]
    3. For each channel c = 1 to 18:
         phi_c(t) = arg(hilbert(x_nb[c,t]))                         [Eq. (3)]
         phi_ref(t) = arg(hilbert(x_ref(t)))                        [Eq. (4)]
         PLV_c = |mean(exp(j*(phi_c - phi_ref)))|                   [Eq. (5)]
         a_c = mean(|hilbert(x_nb[c,t])|)                           [Eq. (6)]
    4. Normalize: a_c = a_c / max(a_c)                               [Eq. (7)]
    5. s_c = PLV_c * a_c                                             [Eq. (8)]
    6. (Optional) Fit sinusoid, compute VE_c                      [Eqs. (9-10)]
    7. Involved = {c : s_c > 0.62}                                  [Eq. (11)]
    8. SE = |Involved| / 18                                          [Eq. (12)]
    return SE, {s_c}, {PLV_c}, {VE_c}
\end{verbatim}

\subsubsection{Output}

\begin{itemize}
  \item \textbf{Spatial extent}: fraction of channels involved, $\mathrm{SE} \in [0, 1]$
  \item \textbf{Per-channel scores}: $\{s_c\}$, $\{\mathrm{PLV}_c\}$, $\{a_c\}$, $\{\mathrm{VE}_c\}$
\end{itemize}

\subsubsection{Performance}

The threshold $\theta_s = 0.62$ was calibrated to maximize agreement with expert spatial extent ratings. The automated spatial extent estimates achieve an intraclass correlation coefficient (ICC) of 0.371 against expert consensus, compared to an inter-rater ICC of 0.373 between independent expert raters. This near-parity with human agreement indicates that the PLV-amplitude method captures the same information experts use when judging RDA spatial extent, while acknowledging that spatial extent is inherently a low-reliability measure even among trained electroencephalographers.

\subsection{BIPD detection: bilateral independent periodic discharges}
\label{app:bipd}

\subsubsection{Overview}

Bilateral independent periodic discharges (BIPD) are defined by the presence of independent periodic discharge trains on each hemisphere, with distinct timing sequences that are not phase-locked to one another. We detect BIPD using a two-stage approach: (1) independent per-hemisphere discharge detection using the HemiCET+DP pipeline, and (2) gradient-boosted tree (GBT) classification of the resulting bilateral timing sequences. The classifier distinguishes BIPD (independent bilateral timing) from GPD (synchronous bilateral discharges) based on features capturing frequency asymmetry, temporal independence, and phase consistency.

\subsubsection{Input}

\begin{itemize}
  \item 18-channel bipolar EEG partitioned into left-hemisphere channels $\mathcal{L}$ ($|\mathcal{L}| = 8$, plus 2 midline) and right-hemisphere channels $\mathcal{R}$ ($|\mathcal{R}| = 8$, plus 2 midline)
\end{itemize}

\subsubsection{Mathematical formulation}

\paragraph{Stage 1: per-hemisphere discharge detection.}
Run the HemiCET+DP pipeline independently on each hemisphere, yielding two discharge time sequences:
\[ \text{Left:}\quad \{t_1^L, t_2^L, \ldots, t_m^L\} \tag{1} \]
\[ \text{Right:}\quad \{t_1^R, t_2^R, \ldots, t_n^R\} \tag{2} \]
The inter-pulse intervals (IPIs) for each hemisphere are:
\[ \mathrm{IPI}_i^L = t_{i+1}^L - t_i^L, \quad i = 1,\ldots,m-1 \tag{3} \]
\[ \mathrm{IPI}_j^R = t_{j+1}^R - t_j^R, \quad j = 1,\ldots,n-1 \tag{4} \]

\paragraph{Stage 2: feature extraction.}
The following features are computed from the bilateral timing sequences.

\paragraph{Feature 1: frequency ratio.}
Per-hemisphere frequency is estimated from the median IPI:
\[ f_L = \frac{1}{\mathrm{median}(\mathrm{IPI}^L)}, \qquad f_R = \frac{1}{\mathrm{median}(\mathrm{IPI}^R)} \tag{5} \]
\[ r_f = \frac{f_L}{f_R} \tag{6} \]
For BIPD, the frequency ratio typically deviates from unity since the two hemispheres generate discharges at independent rates.

\paragraph{Feature 2: regularity (coefficient of variation).}
\[ \mathrm{CoV}_L = \frac{\mathrm{std}(\mathrm{IPI}^L)}{\mathrm{mean}(\mathrm{IPI}^L)}, \qquad \mathrm{CoV}_R = \frac{\mathrm{std}(\mathrm{IPI}^R)}{\mathrm{mean}(\mathrm{IPI}^R)} \tag{7} \]
Low CoV indicates regular periodicity; BIPD hemispheres may differ in regularity.

\paragraph{Feature 3: matched fraction.}
Quantifies the temporal overlap between left and right discharge times using a 100 ms coincidence window:
\[ \mathrm{MF} = \frac{\left|\left\{(i,j) : |t_i^L - t_j^R| < 100\,\mathrm{ms}\right\}\right|}{\max(m, n)} \tag{8} \]
For GPD (synchronous), $\mathrm{MF} \approx 1$; for BIPD (independent), $\mathrm{MF}$ is low and determined by chance coincidence.

\paragraph{Feature 4: phase consistency.}
Compute the nearest-neighbor phase for each left discharge relative to the right hemisphere cycle:
\[ \psi_i = 2\pi \cdot \frac{t_i^L - t_{j^*(i)}^R}{\mathrm{IPI}_{j^*(i)}^R} \tag{9} \]
where $j^*(i) = \operatorname{argmin}_j |t_i^L - t_j^R|$ subject to $t_j^R \leq t_i^L$. Phase consistency is quantified by the mean resultant length:
\[ \mathrm{PC} = \left|\frac{1}{m} \sum_{i=1}^{m} \exp(j \psi_i)\right| \tag{10} \]
For GPD, $\mathrm{PC} \approx 1$ (fixed phase relationship); for BIPD, $\mathrm{PC} \approx 0$ (uniformly distributed phases).

\paragraph{Features 5--7: asymmetry measures.}
\[ \text{Sequence length ratio:} \quad r_n = \frac{\min(m,n)}{\max(m,n)} \tag{11} \]
\[ \text{Frequency asymmetry:} \quad A_f = \frac{|f_L - f_R|}{f_L + f_R} \tag{12} \]
\[ \text{Count asymmetry:} \quad A_n = \frac{|m - n|}{m + n} \tag{13} \]

\paragraph{Stage 3: GBT classification.}
The feature vector $\mathbf{x} = [r_f,\; \mathrm{CoV}_L,\; \mathrm{CoV}_R,\; \mathrm{MF},\; \mathrm{PC},\; r_n,\; A_f,\; A_n]^{\!\top} \in \mathbb{R}^8$ is input to a gradient-boosted tree classifier:
\[ P(\text{BIPD} \mid \mathbf{x}) = \sigma\!\left(\sum_{m=1}^{M} \eta \cdot h_m(\mathbf{x})\right) \tag{14} \]
where $\sigma$ is the logistic function, $h_m$ are regression trees, $\eta$ is the learning rate, and $M$ is the number of boosting rounds.

\paragraph{Training data generation.}
The GBT is trained on synthetic data constructed from real single-hemisphere discharge sequences:
\begin{itemize}
  \item \textbf{Synthetic BIPD}: Draw one LPD sequence from Patient A (left) and an independent LPD sequence from Patient B (right). The cross-patient pairing ensures genuinely independent timing.
  \item \textbf{Synthetic GPD}: Take a single LPD sequence and create a bilateral version by phase-shifting: $t_j^R = t_j^L + \delta$ where $\delta \sim \mathcal{U}(-25\,\mathrm{ms},\; 25\,\mathrm{ms})$. This produces synchronized bilateral discharges with slight jitter.
\end{itemize}

\subsubsection{Algorithm}

\begin{verbatim}
BIPD_DETECT(x_left[c,t], x_right[c,t]):
    # --- Stage 1: Per-hemisphere detection ---
    1. {t^L_1,...,t^L_m} = HemiCET_DP(x_left)        [Eq. (1)]
    2. {t^R_1,...,t^R_n} = HemiCET_DP(x_right)        [Eq. (2)]

    # --- Stage 2: Feature extraction ---
    3. IPI^L, IPI^R from discharge times               [Eqs. (3-4)]
    4. f_L, f_R = 1/median(IPI^L), 1/median(IPI^R)    [Eq. (5)]
    5. r_f = f_L / f_R                                 [Eq. (6)]
    6. CoV_L, CoV_R = std/mean of IPIs                 [Eq. (7)]
    7. MF = matched fraction with 100ms window          [Eq. (8)]
    8. PC = phase consistency (mean resultant length)    [Eqs. (9-10)]
    9. r_n, A_f, A_n = asymmetry features               [Eqs. (11-13)]

    # --- Stage 3: Classification ---
   10. x = [r_f, CoV_L, CoV_R, MF, PC, r_n, A_f, A_n]
   11. P(BIPD) = GBT.predict_proba(x)                   [Eq. (14)]

    return P(BIPD), {t^L}, {t^R}
\end{verbatim}

\subsubsection{Output}

\begin{itemize}
  \item $P(\text{BIPD})$: probability of bilateral independent periodic discharges
  \item Per-hemisphere discharge times for downstream analysis
\end{itemize}

\subsubsection{Performance}

The BIPD classifier achieves an AUC of 0.840 on held-out data, with 63\% sensitivity at the operating threshold. From 198 candidate segments flagged by the detector, 21 were confirmed as true BIPD by expert review. The most discriminative features are matched fraction ($\mathrm{MF}$) and phase consistency ($\mathrm{PC}$), which together capture the defining characteristic of BIPD: temporal independence between hemispheres. The synthetic training strategy avoids the need for large annotated BIPD datasets, which are rare in clinical practice.
